\section{Check the complete toy calculation}
Represent decimal 65, 17 and 3233 in separate sixteen-byte objects, with unused high
limbs zero. Call \texttt{rsa\_pow128} with those base, exponent and modulus pointers.
Inspect status zero and all sixteen output bytes. The result is decimal 2790,
hexadecimal \texttt{0ae6}; its first two bytes are \texttt{e6 0a}, followed by fourteen
zero bytes. Inspecting only the first word can hide a stale high limb.

Copy the resulting numerical ciphertext into the next call's base object, use exponent
2753, and keep modulus 3233. The result must be 65. Perform both calls in the same
stopped execution session so the second call exercises ordinary reuse of stack and
temporary storage. Do not repair the inputs between calls by silently substituting
a known ciphertext constant.

Then vary a boundary at a time. Try exponent zero, exponent one, base zero, a base
larger than its modulus, and modulus one. Set only exponent bit 127, and then set all
128 exponent bits. A function that reads only the first exponent word can pass the
toy calculation and fail both larger cases. Fill scratch storage with a nonzero pattern
before calling to reveal an accumulator or carry that was never initialized.

Finally, put the same objects above address \(2^{32}\), with a code return address
above that boundary too. Their arithmetic values have not changed. A failure there
points toward pointer truncation, a narrow saved address or return-state handling,
not a different mathematical modulus. Caller objects at an offset of four bytes test
the promised four-byte alignment rather than accidental eight-byte alignment.
